% Options for packages loaded elsewhere
\PassOptionsToPackage{unicode}{hyperref}
\PassOptionsToPackage{hyphens}{url}
\documentclass[
  11pt,
]{article}
\usepackage{xcolor}
\usepackage[margin=1in]{geometry}
\usepackage{amsmath,amssymb}
\setcounter{secnumdepth}{-\maxdimen} % remove section numbering
\usepackage{iftex}
\ifPDFTeX
  \usepackage[T1]{fontenc}
  \usepackage[utf8]{inputenc}
  \usepackage{textcomp} % provide euro and other symbols
\else % if luatex or xetex
  \usepackage{unicode-math} % this also loads fontspec
  \defaultfontfeatures{Scale=MatchLowercase}
  \defaultfontfeatures[\rmfamily]{Ligatures=TeX,Scale=1}
\fi
\usepackage{lmodern}
\ifPDFTeX\else
  % xetex/luatex font selection
  \setmainfont[]{Times New Roman}
\fi
% Use upquote if available, for straight quotes in verbatim environments
\IfFileExists{upquote.sty}{\usepackage{upquote}}{}
\IfFileExists{microtype.sty}{% use microtype if available
  \usepackage[]{microtype}
  \UseMicrotypeSet[protrusion]{basicmath} % disable protrusion for tt fonts
}{}
\makeatletter
\@ifundefined{KOMAClassName}{% if non-KOMA class
  \IfFileExists{parskip.sty}{%
    \usepackage{parskip}
  }{% else
    \setlength{\parindent}{0pt}
    \setlength{\parskip}{6pt plus 2pt minus 1pt}}
}{% if KOMA class
  \KOMAoptions{parskip=half}}
\makeatother
\setlength{\emergencystretch}{3em} % prevent overfull lines
\providecommand{\tightlist}{%
  \setlength{\itemsep}{0pt}\setlength{\parskip}{0pt}}
\usepackage{bookmark}
\IfFileExists{xurl.sty}{\usepackage{xurl}}{} % add URL line breaks if available
\urlstyle{same}
\hypersetup{
  hidelinks,
  pdfcreator={LaTeX via pandoc}}

\author{}
\date{}

\begin{document}

\section{Response to the Academic
Editor}\label{response-to-the-academic-editor}

\textbf{Manuscript:} \emph{Proxy-Based Diagnostics of Quantum Oracle
Sketching Robustness for Non-IID Sensor and Telemetry Streams}\\
\textbf{Journal:} \emph{Sensors} (MDPI) · \textbf{Manuscript ID:}
sensors-4470240

\textbf{Editor's comment.} \emph{``I strongly invite the Authors to not
include preprints among references. These works are still undergoing the
peer-review process, thus meaning that their content may change prior to
publication (if they will be finally published). The Authors can
substitute them with related works already published.''}

\textbf{Response.} We thank the Academic Editor for this helpful
recommendation. In accordance with the Editor's request, we have removed
the Zhao et al.~(2026) preprint, ``Exponential quantum advantage in
processing massive classical data'' (arXiv:2604.07639), from the revised
manuscript and from the reference list. The revised bibliography
contains 54 references and no preprints.

Before editing, we audited every in-text use of the preprint together
with the directly dependent clauses. We also examined the published
works already cited in the manuscript --- Kallaugher (FOCS 2021),
Kallaugher, Parekh and Voronova (STOC 2024 and SOSA 2025), and Huang et
al.~(\emph{Science} 2022). These establish important but different
results: problem-specific quantum-streaming space advantages, a
black-box quantum-sketch construction, and quantum advantage for
learning properties of quantum systems. None of them establishes the
framework-specific classification, dimension-reduction, correlated-data,
repetition-overhead, refreshing-time, or classical-hardness claims for
which the preprint had been cited. We therefore did not substitute any
of them for it, since doing so would attribute results to authors who
did not obtain them.

Instead, only the affected statements were revised. Framework-specific
theorem and provenance claims were removed; each remaining reference is
now cited only for the result it actually establishes; and the
operational quantities τ, r, T\_eff, ε and the accuracy proxy are stated
unambiguously as constructs of the present paper. Section 3.3
accordingly separates published theory and its boundary, this paper's
assumptions, its heuristic constructions, its illustrative memory
references, and its scope limitations, and states explicitly that the
published quantum-streaming results do not establish the operational
model used here and that this paper adds no quantum theorem. Table 1's
provenance row and the associated memory-reference descriptions are
relabelled as paper-defined illustrative scaling references. The
remaining Zhao-dependent clauses in the abstract and introduction were
minimally corrected.

The reference audit that accompanied this revision also confirmed the
published form of the Kallaugher, Parekh and Voronova work (SOSA 2025,
pp.~9--45, DOI 10.1137/1.9781611978315.2), added the published STOC 2024
result (DOI 10.1145/3618260.3649709) for the exponential
quantum--classical space separation --- which corrected a Related Work
sentence that would otherwise have credited the 2021 FOCS paper with an
exponential separation, whereas the advantage established there is
polynomial --- and corrected a transposed article number in the Su, Guo
and Wang (2024) entry, now \emph{Information Sciences} \textbf{676},
article 120799, DOI 10.1016/j.ins.2024.120799.

No replacement reference was added to preserve the citation count, and
no experimental data, methodology, equations, proxy definitions,
datasets, seeds, numerical results, figure artwork or conclusions were
changed. All reviewer-requested revisions remain in place, every in-text
citation was re-checked against the bibliography, and the manuscript
compiles with no unresolved citations or references.

On behalf of all authors,

\textbf{AbdelMoniem Helmy} (corresponding author)
\href{mailto:abdelmoniem.hafez@cu.edu.eg}{\nolinkurl{abdelmoniem.hafez@cu.edu.eg}}
· ORCID 0000-0001-5996-6019

\end{document}
